\documentclass{article}

\usepackage[greek,english]{babel}
\usepackage[utf8x]{inputenc}
\usepackage{natbib}
\usepackage{graphicx}
\usepackage{subfigure}
\graphicspath{ {.//} }
\usepackage{xcolor}
\graphicspath{ {./images/} }
\usepackage{titling}

\setlength{\droptitle}{10em}

\title{%
  \selectlanguage{greek}\Huge
 Μοντέλα υπολογιστικού νέφους \\}
  


\begin{document}
\selectlanguage{greek}

\author{\selectlanguage{greek}\LARGE
  Γεώργιος Νικολάου\\
   \texttt{\selectlanguage{english}\large sdi1800134}
  \and
  \selectlanguage{greek}\LARGE
  Νεφέλη Ταβουλάρη\\
   \texttt{\selectlanguage{english}\large sdi1800190}
}

\maketitle


\begin{figure}
\centering
\begin{subfigure}
  \centering
  \includegraphics[width=40mm]{dit_logo}
  \label{fig:sub1}
\end{subfigure}%
\begin{subfigure}
  \centering
  \includegraphics[width=50mm]{NKUA_logo}
  \label{fig:sub2}
\end{subfigure}
\label{fig:test}
\end{figure}




\newpage
\tableofcontents
\newpage
\section{Εισαγωγή}
\section{Σύγχρονα Πληροφοριακά Συστήματα}
\subsection{Ανάπτυξη}
\subsection{Νέες Ανάγκες}
\section{Υπολογιστικό Νέφος}
\subsection{Οφέλη}
\subsection{Προκλήσεις}
\section{Μορφές υπολογιστικού νέφους}
\newpage
\subsection{Λογισμικό ως Υπηρεσία}
Αναμφισβήτητα,η αγορά λογισμικού για τις επιχειρηματικές ανάγκες μιας εταιρείας αποτελεί μια κρίσιμη και περίπλοκη διαδικασία. Στόχος είναι πάντα η ελαχιστοποίηση του κόστους , αλλά και η αύξηση της παραγωγικότητας, επενδύοντας σε όσο μεγαλύτερους αυτοματισμούς. Συνηθισμένα παραδείγματα λογισμικού που χρησιμοποιούνται από τις επιχειρήσεις είναι τα εξής:
\begin{itemize}
  \item Λογισμικό εσωτερικής επικοινωνίας, πχ \selectlanguage{english} Microsoft Teams, Zoom, Webex, Outlook, Jira, Confluence, Slack\selectlanguage{greek}
  \item Λογισμικό για διαχείριση λογιστικών, φορολογικών διαδικασιών
  \item \selectlanguage{english} Enterprise Resource Planning programs (ERP)
  \item Office 365, \selectlanguage{greek}πχ \selectlanguage{english}Word, Excel, Access, Google Docs
  \item Portal \selectlanguage{greek} εταιρείας, διαχείριση αδειών, διαχείριση tickets στο \selectlanguage{english}ΙΤ \selectlanguage{greek}τμήμα \selectlanguage{english}
  \item Version Control System,\selectlanguage{greek} πχ \selectlanguage{english}Github, Gitlab, Bitbucket \selectlanguage{greek}
  \item Λογισμικό διαχείρισης Ανθρωπίνου Δυναμικού \selectlanguage{english}
  \item Customer Relationship Management (CRP), Marketing programs
  \item CAD \selectlanguage{greek}Λογισμικό\selectlanguage{english}
  \item Code Editor,\selectlanguage{greek} πχ\selectlanguage{english} VSCode, IntelliJ, PyCharm\selectlanguage{greek}

\end{itemize}
Παραδοσιακά, αυτή η ανάγκη των εταιρειών ικανοποιούνταν με το λεγόμενο, Λογισμικό ως Προϊόν \selectlanguage{english}(Software as a Product, SaaP), \selectlanguage{greek}βάσει του οποίου, η εταιρεία αγοράζει άδεια για να χρησιμοποιεί τη συγκεκριμένη υπηρεσία, την οποία φιλοξενεί ύστερα, στις εγκαταστάσεις της \selectlanguage{english}(on premises).\selectlanguage{greek} Επίσης, αναλαμβάνει τα έξοδα συντήρησης και ανανέωσης του λογισμικού. \\ \\
Ωστόσο, με την ανάπτυξη του Υπολογιστικού Νέφους, το Λογισμικό ως Υπηρεσία γίνεται όλο και πιο διαδεδομένο.Το Λογισμικό ως Υπηρεσία, γνωστό και ως \selectlanguage{english}Software as a Service (SaaS) \selectlanguage{greek}αναφέρεται στην παροχή ενός λογισμικού / πληροφοριακού συστήματος, το οποίο φιλοξενείται κεντρικά, στο Νέφος. Η φιλοσοφία πίσω από αυτή την διαδεδομένη μέθοδο, είναι πως ο χρήστης μπορεί να απολαμβάνει μια ηλεκτρονική υπηρεσία, χωρίς να χρειάζεται να διαχειρίζεται τις υποδομές, το δίκτυο, τις βάσεις δεδομένων, το ενδιάμεσο λογισμικό, τους διακομιστές / σέρβερς, τα δεδομένα, το λειτουργικό σύστημα, αλλά ούτε και τις εφαρμογές. Εν ολίγοις, ο χρήστης χρησιμοποιεί το Πληροφοριακό Σύστημα , ως συνδρομητής, για να ικανοποιήσει τις ανάγκες του και ο Πάροχος της υπηρεσίας αναλαμβάνει να διαχειρίζεται όλους τους απαιτούμενους πόρους.  \\ 
\newpageΤα κυριότερα πλεονεκτήματα ενός τέτοιου μοντέλου είναι τα εξής :
\begin{itemize}
    \itemΧαμηλότερο κόστος για τις επιχειρήσεις και τους ιδιώτες, οι οποίοι δε χρειάζεται να αγοράσουν το λογισμικό και τις απαραίτητες άδειες. Ανάπτυξη οικονομίας κλίμακας με μείωση κόστους συντήρησης και παροχής του λογισμικού. Επίσης, το κόστος είναι συγκεκριμένο \selectlanguage{english}(pay-as-you-go) \selectlanguage{greek}και πολλές φορές δεν υπάρχει από την αρχή της χρήσης (πχ \selectlanguage{english}Netflix).\selectlanguage{greek} 
    \itemΔεν υπάρχει ανάγκη γνώσης χρήσης των διαφόρων λογισμικών, ούτε επομένως και εξειδικευμένου προσωπικού σε μια επιχείρηση, με αποτέλεσμα ξανά τη μείωση κόστους. Ευχρηστία, σύνδεση στην εφαρμογή μέσω κάποιου \selectlanguage{english}Dashboard\selectlanguage{greek} ή \selectlanguage{english}API.\selectlanguage{greek}
    \itemΠρόσβαση από παντού, μέσω Διαδικτύου και οποιασδήποτε συσκευής, χωρίς να χρειάζεται να εγκατασταθεί πρόσθετο λογισμικό (πχ \selectlanguage{english}dependencies,\selectlanguage{greek} εργαλεία, packages, λειτουργικό σύστημα), ή σύνδεση με \selectlanguage{english}VPN\selectlanguage{greek}. Επιτυγχάνεται, λοιπόν, μεγαλύτερη ευελιξία και αύξηση της παραγωγικότητας 
    \itemΌλοι οι χρήστες έχουν την ίδια έκδοση, η οποία αναβαθμίζεται αυτόματα
    \itemΠρόκειται για ολοκληρωμένα και εύχρηστα Πληροφοριακά Συστήματα, ανεξάρτητα από το Λειτουργικό Σύστημα ή το Υλικό του καθενός, ακολουθούν επομένως τις αρχές της εικονικοποίησης (\selectlanguage{english}virtualization)\selectlanguage{greek} και της αφαιρετικότητας (\selectlanguage{english}abstraction)\selectlanguage{greek}
    \itemΑποτελεί μια πρακτική που μπορεί να αποβεί ωφέλιμη και για το περιβάλλον, καθώς χρησιμοποιούνται όσο το δυνατόν λιγότεροι πόροι 
    \itemΚαλύτερη συνεργασία μέσα σε μια εταιρεία, λόγω του \selectlanguage{english}multitenancy,\selectlanguage{greek} δηλαδή της παράλληλης χρήσης της ίδιας εφαρμογής από πολλούς χρήστες. Χαρακτηριστικό του \selectlanguage{english}Multitenancy\selectlanguage{greek} είναι πως οι χρήστες νιώθουν σα να έχουν αποκλειστική χρήση του Λογισμικού
    \itemΑσφάλεια από κακόβουλους εισβολείς, προστασία, απόδοση, προσβασιμότητα, υψηλή διαθεσιμότητα, κλιμάκωση, αξιοπιστία και γενικότερα, όλα τα γνωστά πλεονεκτήματα του\selectlanguage{english} Cloud Computing\selectlanguage{greek}. Καθώς οι διάφορες εφαρμογές φιλοξενούνται σε κορυφαία \selectlanguage{english}Data Centers,\selectlanguage{greek} ειδικά σχεδιασμένα με αυτό το σκοπό, σε υψηλού επιπέδου υποδομές, με \selectlanguage{english}firewalls,\selectlanguage{greek} κρυπτογράφηση και ελεγχόμενη πρόσβαση
    \itemΧρησιμοποιώντας \selectlanguage{english}SaaS, \selectlanguage{greek}μια επιχείρηση έχει την ευκαιρία να αφοσιωθεί σε αυτό που την ενδιαφέρει πιο πολύ, το προϊόν της. Το ίδιο βέβαια ισχύει και για τους ιδιώτες. Αύξηση, λοιπόν, της καινοτομίας και της ευκινησίας \selectlanguage{english}(agility)
\selectlanguage{greek}
\end{itemize}
\newpage
Τα κυριότερα μειονεκτήματα :
\begin{itemize}
 \itemΑναγκαία η χρήση Διαδικτύου, διαφορετικά δεν μπορεί να υπάρξει πρόσβαση στην εφαρμογή
\itemΧαμηλή ασφάλεια ευαίσθητων δεδομένων, αν δεν είναι έμπιστος ο πάροχος
\itemΈλλειψη ελέγχου
\itemΤαχύτητα μετάδοσης
\end{itemize}
Συμπέρασμα : \\ \\
Με τη σωστή επιλογή ενός γνωστού και αποδεκτού παρόχου υπηρεσιών λογισμικού, την υπογραφή Συμφωνητικού για το Επίπεδο Λειτουργίας Υπηρεσίας,\selectlanguage{english} SLA (service level agreement),\selectlanguage{greek} όπου καταγράφονται όλες οι νομικές υποχρεώσεις του παρόχου, καθώς και με τη χρήση\selectlanguage{english}  5G,\selectlanguage{greek} το Λογισμικό ως Υπηρεσία είναι μια φτηνή και σίγουρη επιλογή για μια επιχείρηση, ειδικά αν δεν αναζητά κάποια εξειδικευμένη υπηρεσία, όπου έχει ανάγκη\selectlanguage{english}  customization.\selectlanguage{greek}


\newpage
\subsection{Πλατφόρμα ως Υπηρεσία}
Το μοντέλο Πλατφόρμα ως Υπηρεσία \selectlanguage{english}(PaaS)\selectlanguage{greek} επιτρέπει στους χρήστες να επωφελούνται από τις υποδομές, το δίκτυο, το λειτουργικό σύστημα, τις βάσεις δεδομένων, τα οποία διαχειρίζεται και συντηρεί ο πάροχος της υπηρεσίας και ταυτόχρονα να έχουν την ευελιξία να σχεδιάσουν και να υλοποιήσουν ο,τιδήποτε επιθυμούν για την επιχείρησή τους, χρησιμοποιώντας έτοιμες μικρο-υπηρεσίες που υπάρχουν ήδη στο περιβάλλον της πλατφόρμας. Διατίθεται μέσω \selectlanguage{english} public, private \selectlanguage{greek}και \selectlanguage{english}hybrid\selectlanguage{greek} νέφους. \\
Παραδείγματα της Πλατφόρμας ως Υπηρεσία :
\begin{itemize}
\item Πλατφόρμα Επικοινωνίας ως Υπηρεσία, \selectlanguage{english}Communications Platform as a Service    (CPaaS) \selectlanguage{greek},η οποία περιλαμβάνει συνδυασμό ήχου, βίντεο, μηνυμάτων, πχ \selectlanguage{english}Skype, WhatsApp
\item Mobile Platform as a Service, \selectlanguage{greek}για την εύκολη ανάπτυξη εφαρμογών \selectlanguage{english}(Integrated Development Environment, IDE)
\item Open PaaS, open source / \selectlanguage{greek}ανοιχτού κώδικα πλατφόρμα, όπου ο χρήστης μπορεί να αναπτύξει εφαρμογές, πολλές φορές με οποιαδήποτε γλώσσα προγραμματισμού ή βάση δεδομένων επιθυμεί πχ \selectlanguage{english}Google App Engine, Microsoft Azure, Heroku\selectlanguage{greek}
\itemΜια άλλη αξιοσημείωτη χρήση \selectlanguage{english}PaaS \selectlanguage{greek}είναι στα \selectlanguage{english}DevOps tools (CI/CD pipelines), Agile Development
\item Integrations PaaS (iPaas) \selectlanguage{greek}για συγχώνευση δεδομένων, υπηρεσιών κτλ
\item \selectlanguage{english}AIPaaS, \selectlanguage{greek}για δημιουργία μοντέλων Μηχανικής Μάθησης με προεκπαιδευμένο/\selectlanguage{english}pre-trained\selectlanguage{greek} μοντέλα, τα οποία διευκολύνουν κατά πολύ τη διαδικασία 
\end{itemize}
Τα κυριότερα πλεονεκτήματα ενός τέτοιου μοντέλου είναι τα εξής :
\begin{itemize}
\item Οι προγραμματιστές μπορούν να αφοσιωθούν στην ανάπτυξη λογισμικού και συγκεκριμένα στη λογική \selectlanguage{english}(business intelligence)\selectlanguage{greek} πίσω από την εκάστοτε εφαρμογή, αντί να χάνουν χρόνο με την υποδομή, το δίκτυο, το λειτουργικό σύστημα, την αποθήκευση, τους διακομιστές. Όλα αυτά τα αναλαμβάνει ο πάροχος. Μείωση, λοιπόν, του \selectlanguage{english}low level programming,\selectlanguage{greek} το οποίο είναι και το πιο απαιτητικό.
\item Απλό μοντέλο, χαμηλό κόστος για ανάπτυξη εφαρμογών ,\selectlanguage{english} pay-as-you-go\selectlanguage{greek}
\item Προσβασιμότητα, υψηλή διαθεσιμότητα, κλιμάκωση
\item Μείωση απαιτούμενου κώδικα, καθώς παρέχονται έτοιμα πολλά χρήσιμα εργαλεία \selectlanguage{english}(pre-coded components, sample code, software tools, kits, libraries), \selectlanguage{greek}αύξηση της παραγωγικότητας 
\item Εύκολη μεταφορά στο Υβριδικό Νέφος
\item Έτοιμο περιβάλλον για \selectlanguage{english}versioning, build (debugger, compiler, code editor), deployment, testing, monitoring,\selectlanguage{greek} διαθέσιμα\selectlanguage{english} plugins\selectlanguage{greek} και \selectlanguage{english}frameworks\selectlanguage{greek} που διευκολύνουν αισθητά τη χρήση της πλατφόρμας. Επιπρόσθετα,  το μοντέλο \selectlanguage{english}PaaS \selectlanguage{greek}παρέχει συχνά εργαλεία για\selectlanguage{english} Data Analytics, \selectlanguage{greek}για δεδομένα που αφορούν το χρήστη ως προς την εφαρμογή.
\item Παροχή ενσωματωμένης βάσης \selectlanguage{english}(Database as a Service, \selectlanguage{greek}πχ \selectlanguage{english}MongoDB Atlas)\selectlanguage{greek},και γενικότερα έτοιμων και εύχρηστων υπηρεσιών
\item Παρέχει τα θετικά του \selectlanguage{english}SaaS, \selectlanguage{greek}όσον αφορά την ανεξαρτησία από τις υποδομές, αλλά προσφέρει ακόμη και επιπλέον ευελιξία, καθώς ο χρήστης μπορεί να δημιουργήσει ο,τιδήποτε επιθυμεί για να καλύψει πλήρως τις ανάγκες του, κάτι που ίσως δε θα μπορούσε να γίνει με ένα\selectlanguage{english} SaaS,\selectlanguage{greek} το οποίο δίνεται έτοιμο 
\item Παρόμοιο μοντέλο με το Συνάρτηση ως Υπηρεσία\selectlanguage{english} (Function as a Service), Serverless as a Service,\selectlanguage{greek} όπου σημασία έχει η λογική και όχι και η υποδομή
\end{itemize}
Τα κυριότερα μειονεκτήματα :
\begin{itemize}
\item Δύσκολη μετακίνηση σε άλλο πάροχο, καθώς είναι πιθανό εκεί να γίνονται με άλλον τρόπο οι διαδικασίες ή να μην παρέχονται οι ίδιες υπηρεσίες/ευκολίες κτλ
\item Εξάρτηση από τον πάροχο, η τιμή μπορεί να αυξηθεί ανά πάσα στιγμή, αν υπάρξει πρόβλημα στις υποδομές, η εφαρμογή θα καταρρεύσει
\item Έλλειψη ελέγχου, όπως και στο μοντέλο Λογισμικό ως Υπηρεσία, καθώς κάποιος άλλος καθορίζει μεγάλο μέρος των \selectlanguage{english}configurations
\end{itemize}
\selectlanguage{greek}
Συμπέρασμα: \\ \\
Με τη σωστή επιλογή της υπηρεσίας, ώστε να καλύπτει τις ανάγκες της επιχείρησης / του χρήστη, αλλά και να ανήκει σε έμπιστο πάροχο, η Πλατφόρμα ως Υπηρεσία μπορεί να βελτιώσει και να διευκολύνει τη δουλειά του Developer, καθώς δε χρειάζεται να ξεκινάει από το μηδέν για κάθε του εργασία.

\subsection{Υποδομή ως Υπηρεσία}
\section{Από το \selectlanguage{english}IaaS \selectlanguage{greek}στο \selectlanguage{english}SaaS}

\subsection{Σύγκριση}
\section{Συμπεράσματα - Επισημάνσεις}
\end{document}
